\section{PDCA：点検した結果を改善へ戻す}
\label{theory:management:pdca}
計画し、実施し、結果を点検し、是正と改善へ反映する順序がPDCAである。目標と評価方法を定め、実行の記録を取り、目標との差と原因を調べて次の計画へつなぐ\parencite{eth54_pdca2017}。

Checkは単に点検作業をしたという事実ではなく、計画に照らして結果と過程を評価する段階である。Actも行動一般ではなく、その結果を踏まえる是正・改善を表す。研究開発の評価は次の研究や資源配分へ役立てる。事前・中間・事後評価の時点と、PDCAの語順は別々に読む。
